% ===========================================================
%  CHAPTER II — STATE OF THE ART
% ===========================================================
\chapter{State of the Art}

% -----------------------------------------------------------
\section{Human Motor Control: Computational Perspectives}

\subsection{The Motor Control Problem}

Motor control refers to the full set of processes by which the CNS plans,
initiates, and regulates voluntary and involuntary movements. From a computational
perspective, it is an optimal control problem under constraints in a very high-
dimensional state space~\parencite{Todorov2004, Scott2004}.

Three major computational hypotheses guide the modeling of motor control:
\begin{enumerate}
  \item \textbf{\textit{Jerk} minimization}~\parencite{Flash1985}: trajectories
        minimize the derivative of acceleration, producing smooth bell-shaped
        velocity profiles characteristic of human pointing movements.
  \item \textbf{End-point variance minimization}~\parencite{Harris1998}: this
        integrates the stochastic nature of the neural signal
        (\textit{signal-dependent noise}) and predicts curved trajectories for
        large-amplitude movements.
  \item \textbf{Stochastic optimal control}~\parencite{Todorov2002}: this unifies
        the previous approaches within a Bayesian framework and underlies the theory
        of Optimal Feedback Control~(OFC). OFC predicts that the CNS corrects only
        errors that interfere with task achievement (the ``minimum intervention
        principle''), thereby minimizing unnecessary co-contraction.
\end{enumerate}

\subsection{Internal Models and Efference Prediction}

The CNS uses \textit{internal models} to predict the sensory consequences of motor
commands and compensate for feedback delays~\parencite{Wolpert1998}. The MOSAIC
model (\textit{Modular Selection and Identification for Control})~\parencite{Wolpert1998}
proposes a modular architecture in which several internal models combine
adaptively. This perspective is relevant to our work: the learned \DRL{} policies can
be interpreted as \textit{implicit inverse models} of musculoskeletal dynamics.

\subsection{Muscle Synergies and Modular Organization}
\label{sec:synergies_related}

A complementary theory, that of muscle synergies~\parencite{dAvella2003,
Ting2007, Bizzi2013}, proposes that the CNS simplifies control by recruiting
pre-structured neural modules, each activating a coherent subset of muscles.
Mathematically, the observed activations $A(t) \in \R^{N}$ can be decomposed into a
small number $k \ll N$ of synergies $W \in \R^{N \times k}$ activated by temporal
coefficients $C(t) \in \R^k$:
\begin{equation}
  A(t) \approx W \cdot C(t), \quad W_{ij} \geq 0, \; C_{ij} \geq 0
  \label{eq:nmf_principle}
\end{equation}
This decomposition, obtained through non-negative matrix factorization~(NMF), is
robustly observed in humans, primates, and cats across a variety of tasks. Testing
whether a \DRL{} agent produces similar synergies therefore constitutes a highly
relevant \textit{biomimetic} validation.

\subsection{Muscle Redundancy and the Inverse Problem}

The human body exhibits significant muscular redundancy: the upper limb alone is
driven by more than forty muscles, whereas the number of joint degrees of freedom is
on the order of ten. This redundancy implies that an infinite number of muscular
activation combinations can produce the same movement, which constitutes the
\textit{indeterminacy problem} of motor control~\parencite{Zajac1989}.

Classical approaches to resolving this indeterminacy include:
\begin{itemize}
  \item Static optimization with muscle-stress
        minimization~\parencite{Crowninshield1981}:
        $\min \sum_i (F_i^M / F_{0,i}^M)^2$.
  \item Forward simulation with dynamic optimization of
        controls~\parencite{Anderson2001}.
  \item EMG-based approaches (\textit{EMG-driven simulation}).
\end{itemize}
These methods do not capture the learning dynamics characteristic of biological motor
control, which motivates the \DRL{} approach.

% -----------------------------------------------------------
\section{Musculoskeletal Modeling}

\subsection{Simulation Platforms}

Several software platforms support musculoskeletal simulation
(Table~\ref{tab:platforms}).

\begin{table}[H]
  \centering
  \small
  \caption{Comparison of the main musculoskeletal simulation platforms.}
  \label{tab:platforms}
  \rowcolors{2}{lightblue!30}{white}
  \begin{tabularx}{\textwidth}{@{} L{2.4cm} L{2.6cm} X c @{}}
    \toprule
    \rowcolor{lightblue}
    \textbf{Platform} & \textbf{License} & \textbf{Strengths} & \textbf{Ref.} \\
    \midrule
    OpenSim     & Open-source & Validated anatomical models, rich GUI, motion library & \cite{Delp2007} \\
    AnyBody     & Commercial  & Static/dynamic analysis, subject-specific parameters & --- \\
    \MuJoCo     & Free         & Speed, robust contact handling, Python interface, native muscle actuators & \cite{Todorov2012} \\
    MyoSuite    & Open-source & Musculoskeletal + DRL, detailed hand ($>$50~muscles) & \cite{Caggiano2022} \\
    dm\_control & Open-source & Standardized RL environments, based on \MuJoCo{} & \cite{Tassa2020} \\
    MotorNet    & Open-source & Differentiable, native PyTorch integration, planar arm & \cite{Codol2023} \\
    \bottomrule
  \end{tabularx}
  \normalsize
\end{table}

\MuJoCo{}~\parencite{Todorov2012}, acquired by DeepMind in 2021 and later made
freely available, represents the state of the art for \DRL{} thanks to its
exceptional performance: up to one million simulation steps per second on a standard
CPU, a robust contact solver, and native support for Hill-type muscle actuators.

\subsection{Computational Bottleneck: the Real-Time Cost of Exact Solving}
\label{sec:compute_bottleneck}

A critical and often overlooked aspect of physics-based muscle control is its
per-step computational cost. Because the fiber--tendon equilibrium
equation~\eqref{eq:equilibrium} is implicit and nonlinear in~$l^M$, it cannot be
solved analytically: Newton--Raphson iteration is required at \emph{every
simulation timestep, for every muscle}. For a 9-muscle model running at
\SI{100}{\hertz}, this amounts to approximately
$9 \times (5\text{--}40) \times 100 = 4\,500$--$36\,000$~arithmetic iterations
per second --- a workload that grows linearly with model complexity and control
frequency.

On a standard CPU, each Newton--Raphson solve requires 5--40 iterations and
takes on the order of 10--100~\si{\micro\second} per muscle. While this is
manageable in an offline simulation context, it represents a non-trivial burden
for embedded systems (\eg{} wearable exoskeletons, myoelectric prostheses) where
power and memory are constrained.

This bottleneck motivates the central hypothesis of the present work: a
\emph{trained RL policy}, which replaces iterative solving with a single matrix
multiplication chain through a shallow MLP, can produce equally accurate muscle
activation commands in \emph{microseconds}, with all the ``computation'' amortized
into a one-time training phase. Whether this hypothesis holds --- and to what extent
the quality of control is preserved --- is one of the core empirical questions
addressed in Chapter~\ref{chap:impl} and the computational benchmark of
Chapter~\ref{chap:drl} (Section~\ref{sec:compute_efficiency}).

\subsection{The Hill Muscle Model: Complete Formulation}
\label{sec:hill}

The Hill muscle model~\parencite{Hill1938}, in its modern
formulation~\parencite{Zajac1989, Millard2013}, is the reference biophysical
representation. It decomposes the muscle--tendon complex into three functional
elements (Figure~\ref{fig:hill}):

\begin{itemize}
  \item \textbf{Contractile element~(CE)}: generates active force and is governed by
        activation $a(t) \in [0,1]$.
  \item \textbf{Series elastic element~(SEE)}: models tendon elasticity,
        storing and returning elastic energy.
  \item \textbf{Parallel elastic element~(PEE)}: models the passive stiffness
        of muscle tissue (fascia, sarcolemma).
\end{itemize}

\begin{figure}[H]
  \centering
  \resizebox{0.96\textwidth}{!}{%
  \begin{tikzpicture}[x=1.15cm, y=0.82cm]
    % PEE
    \draw[thick, accentblue] (0,1.55) -- (0.5,1.55);
    \draw[thick, accentblue, decorate, decoration={zigzag, segment length=4pt, amplitude=3pt}]
         (0.5,1.55) -- (2.5,1.55);
    \draw[thick, accentblue] (2.5,1.55) -- (3,1.55);
    \node[above, accentblue, font=\footnotesize] at (1.5,1.95) {PEE};
    % CE
    \draw[thick, navyblue, fill=lightblue!40] (0.5,0.5) rectangle (2.5,1.15);
    \node[navyblue, font=\footnotesize] at (1.5,0.83)
         {CE \quad $f_\text{FL}(\lnorm)\,f_\text{FV}(\vnorm)$};
    % Outer box
    \draw[thick] (0,0.3) -- (0,1.8);
    \draw[thick] (0,0.3) -- (0.5,0.3) -- (0.5,0.5);
    \draw[thick] (0,1.8) -- (0.5,1.8) -- (0.5,1.55);
    \draw[thick] (2.5,0.5) -- (2.5,1.8) -- (3,1.8);
    \draw[thick] (2.5,0.3) -- (3,0.3);
    \draw[thick] (3,0.3) -- (3,1.8);
    % Pennation marker
    \draw[dashed, midgrey] (3,1.05) -- (3.35,1.05);
    \draw[midgrey, ->, >=stealth] (3.25,1.05) arc[start angle=0, end angle=15, radius=0.3];
    \node[right, midgrey, font=\tiny] at (3.4,1.15) {$\alpha$};
    % SEE
    \draw[thick, successgreen] (3,1.05) -- (3.5,1.05);
    \draw[thick, successgreen, decorate, decoration={zigzag, segment length=4pt, amplitude=3pt}]
         (3.5,1.05) -- (5.0,1.05);
    \draw[thick, successgreen] (5.0,1.05) -- (5.5,1.05);
    \node[above, successgreen, font=\footnotesize] at (4.25,1.4) {SEE (tendon)};
    % Anchors
    \draw[thick,|-] (-0.5,1.05) -- (0,1.05);
    \draw[thick,-|] (5.5,1.05) -- (6.0,1.05);
    \node[font=\footnotesize] at (-1.0,1.05) {Orig.};
    \node[font=\footnotesize] at (6.55,1.05) {Ins.};
    % Lengths
    \draw[<->, thick, midgrey] (0,-0.15) -- (3,-0.15)
          node[midway, below, font=\footnotesize] {$l^M$ (fibre)};
    \draw[<->, thick, successgreen] (3,-0.15) -- (5.5,-0.15)
          node[midway, below, font=\footnotesize] {$l^T$ (tendon)};
  \end{tikzpicture}%
  }
  \caption{Complete diagram of the Hill muscle--tendon model. The pennation
           angle~$\alpha$ between the fiber and the line of action is shown at the
           insertion point. The total length of the muscle--tendon unit is
           $l^\text{MT} = l^M \cos\alpha + l^T$.}
  \label{fig:hill}
\end{figure}

\subsubsection{Force-Length Relationship (FL)}

Maximum isometric force varies as a function of normalized muscle length according
to a bell-shaped curve, reflecting the optimal \textit{overlap} of actin and myosin
filaments~\parencite{Gordon1966}:
\begin{equation}
  f_\text{FL}(\lnorm) = \exp\!\left(-\left(\frac{\lnorm - 1}{\gamma_\text{FL}}\right)^{\!2}\right),
  \qquad \lnorm = \frac{l^M}{l^M_\text{opt}}
  \label{eq:fl}
\end{equation}
where $\gamma_\text{FL} \approx 0{,}45$ is a shape parameter and $l^M_\text{opt}$ is
the optimal fiber length at which the muscle produces its maximum isometric
force~$F_0^M$.

\subsubsection{Force-Velocity Relationship (FV)}

Hill's equation~\parencite{Hill1938} describes the hyperbolic force--velocity
relationship for concentric contractions; a linear extension is used for eccentric
contractions~\parencite{Zajac1989}:
\begin{equation}
  f_\text{FV}(\vnorm) =
  \begin{cases}
    \dfrac{v_\text{max} - \vnorm}{v_\text{max} + k_\text{ce}\,\vnorm}
      & \vnorm \leq 0 \quad \text{(concentric)} \\[10pt]
    \dfrac{f_\text{max} + k_\text{ec}\,\vnorm}{1 + k_\text{ec}\,\vnorm}
      & \vnorm > 0 \quad \text{(eccentric)}
  \end{cases}
  \label{eq:fv}
\end{equation}
where $\vnorm = \dot{l}^M / (v_\text{max}\,l^M_\text{opt})$ is the normalized
contraction velocity, $v_\text{max} \approx 10$ (optimal lengths/s),
$k_\text{ce} \approx 0{,}25$, and $f_\text{max} \approx 1{,}3$ (maximum eccentric
coefficient).

\subsubsection{Parallel Elastic Element (PEE)}

The passive force generated by muscle tissue when
$\lnorm > 1$~\parencite{Zajac1989}:
\begin{equation}
  F_\text{PEE}(\lnorm) =
  \begin{cases}
    F_0^M \cdot k_\text{PE} \left(\lnorm - 1\right)^2 & \lnorm > 1 \\
    0 & \lnorm \leq 1
  \end{cases}
  \label{eq:pee}
\end{equation}
with $k_\text{PE} \approx 5{,}0$ a dimensionless stiffness coefficient such that
$F_\text{PEE}(\lnorm = 1{,}6) \approx F_0^M$.

\subsubsection{Series Elastic Element (SEE) --- Tendon}

The tendon is modeled by a quadratic law for small strains and a quasi-linear one
beyond the zero-load length
$l^T_\text{slack}$~\parencite{Thelen2003}:
\begin{equation}
  F_\text{SEE}(\lnorm^T) =
  \begin{cases}
    0 & \lnorm^T \leq 1 \\
    F_0^M \dfrac{(\lnorm^T - 1)^2}{k_T + \lnorm^T - 1} & \lnorm^T > 1
  \end{cases}
  \label{eq:see}
\end{equation}
where $\lnorm^T = l^T / l^T_\text{slack}$ is the normalized tendon length and
$k_T \approx 0{,}0475$ is a dimensionless stiffness coefficient.

\subsubsection{Pennation Angle Dynamics}

The pennation angle $\alpha$ (angle between the muscle fibers and the tendon line of
action) varies with fiber length according to the constant-volume
constraint~\parencite{Zajac1989}:
\begin{equation}
  \alpha(l^M) = \arcsin\!\left(\frac{\sin\alpha_0 \cdot l^M_\text{opt}}{l^M}\right)
  \label{eq:pennation}
\end{equation}
where $\alpha_0$ is the pennation angle at optimal length. The muscle force
transmitted to the tendon is projected by $\cos\alpha$. The total length of the
muscle--tendon unit satisfies the geometric constraint:
\begin{equation}
  l^\text{MT} = l^M \cos\alpha(l^M) + l^T
  \label{eq:lmt}
\end{equation}

\subsubsection{Activation Dynamics}

Muscle activation dynamics introduce an electromechanical delay between the neural
command $u(t)$ and the activation $a(t)$~\parencite{Thelen2003}:
\begin{equation}
  \dot{a}(t) = \frac{u(t) - a(t)}{\tau(u,a)}, \qquad
  \tau(u,a) = \begin{cases} \tau_\text{act}   & u(t) > a(t) \\
                            \tau_\text{deact} & u(t) \leq a(t) \end{cases}
  \label{eq:activation}
\end{equation}
where $u(t) \in [0,1]$ is the neural excitation, $a(t) \in [0,1]$ the activation,
and $\tau_\text{act} \approx 10$--$20$~ms, $\tau_\text{deact} \approx 40$--$70$~ms.

\subsubsection{Fiber-Tendon Equilibrium Equation and Total Force}

The quasi-static equilibrium between contractile-element force, passive force,
and tendon force gives:
\begin{equation}
  F_\text{SEE}(l^T) = \left[a\,F_0^M\,f_\text{FL}(\lnorm)\,f_\text{FV}(\vnorm)
                      + F_\text{PEE}(\lnorm)\right] \cos\alpha
  \label{eq:equilibrium}
\end{equation}

Given $l^\text{MT}$ (computed from kinematics),
Equation~\eqref{eq:equilibrium} combined with the geometric
constraint~\eqref{eq:lmt} defines an implicit nonlinear equation in $l^M$
(fiber length), solved numerically at each simulation step
(Appendix~\ref{app:newton}).

The total muscle force applied to the skeleton is:
\begin{equation}
  F^\text{muscle}(a,l^M,\dot{l}^M) = F_\text{SEE}(l^T)
  \label{eq:fmuscle}
\end{equation}

% -----------------------------------------------------------
\section{Deep Reinforcement Learning for Continuous Control}

\subsection{Theoretical Foundations}
\label{sec:rl_foundations}

Reinforcement learning~(\RL{}) is formalized as a Markov Decision
Process~(MDP) defined by the five-tuple
$(\mathcal{S}, \mathcal{A}, P, R, \gamma)$~\parencite{Bellman1957, Sutton2018}:
\begin{itemize}
  \item $\mathcal{S}$: state space (continuous, $\subset \R^{38}$ in this work).
  \item $\mathcal{A}$: action space (continuous, $= [0,1]^9$ in this work).
  \item $P : \mathcal{S} \times \mathcal{A} \times \mathcal{S} \to [0,1]$:
        transition probability.
  \item $R : \mathcal{S} \times \mathcal{A} \to \R$: reward signal.
  \item $\gamma \in [0,1)$: temporal discount factor.
\end{itemize}

The objective is to find a policy
$\pi : \mathcal{S} \to \mathcal{A}$ that maximizes the expected discounted return:
\begin{equation}
  J(\pi) = \E_{\tau \sim \pi}\!\left[\sum_{t=0}^{\infty} \gamma^t\,r(s_t, a_t)\right]
  \label{eq:rl_objective}
\end{equation}

The modern actor--critic framework~\parencite{Konda2000} combines value estimation
(the critic) and policy improvement (the actor), and forms the basis of DDPG, TD3,
and SAC.

\subsection{Policy Gradient Theorem}

The policy-gradient theorem~\parencite{Sutton2018} provides the gradient of
$J(\pi_\theta)$ with respect to the parameters $\theta$:
\begin{equation}
  \nabla_\theta J(\pi_\theta) = \E_{\tau \sim \pi_\theta}\!\left[
    \sum_{t=0}^{T} \nabla_\theta \log \pi_\theta(a_t|s_t) \cdot G_t
  \right]
  \label{eq:pg_theorem}
\end{equation}
where $G_t = \sum_{k=t}^{T} \gamma^{k-t}\,r_k$ is the cumulative return from step~$t$.
The variance of this estimator is reduced by introducing an advantage function
$A^\pi(s,a) = Q^\pi(s,a) - V^\pi(s)$ in place of~$G_t$.

\subsection{State-of-the-Art Algorithms}
\label{sec:algorithms}

\subsubsection{DDPG — Deep Deterministic Policy Gradient}

\textcite{Lillicrap2016} introduce DDPG, the first modern actor-critic algorithm
for continuous action spaces. The actor learns a deterministic policy
$\mu_\theta : \mathcal{S} \to \mathcal{A}$; the critic evaluates $Q_\phi(s,a)$.

Critic update (Bellman-error minimization):
\begin{equation}
  \mathcal{L}(\phi) = \E_{(s,a,r,s') \sim \mathcal{B}}\!\left[\bigl(Q_\phi(s,a) - y\bigr)^2\right],
  \quad y = r + \gamma\,Q_{\phi'}(s', \mu_{\theta'}(s'))
  \label{eq:ddpg_critic}
\end{equation}

Actor update:
\begin{equation}
  \nabla_\theta J \approx \E\!\left[\nabla_a Q_\phi(s,a)\big|_{a=\mu_\theta(s)}
  \cdot \nabla_\theta \mu_\theta(s)\right]
  \label{eq:ddpg_actor}
\end{equation}

Target-network update (\textit{Polyak averaging}, $\tau = 0{,}005$):
\begin{equation}
  \phi' \leftarrow \tau\phi + (1-\tau)\phi', \qquad
  \theta' \leftarrow \tau\theta + (1-\tau)\theta'
  \label{eq:polyak}
\end{equation}

Exploration is ensured by an additive Ornstein--Uhlenbeck process on the
actions. DDPG suffers from systematic overestimation of $Q$ values and strong
sensitivity to hyperparameters~\parencite{Fujimoto2018}.

\subsubsection{TD3 — Twin Delayed DDPG}

\textcite{Fujimoto2018} correct DDPG's $Q$-overestimation bias through three
complementary mechanisms:

\begin{enumerate}
  \item \textbf{Twin critics}: two critics $Q_{\phi_1}, Q_{\phi_2}$; the target
        uses the minimum:
        \begin{equation}
          y = r + \gamma \min_{i=1,2} Q_{\phi'_i}\!\left(s',\, \tilde{a}\right),
          \quad \tilde{a} = \mu_{\theta'}(s') + \epsilon, \;
          \epsilon \sim \mathrm{clip}\!\left(\mathcal{N}(0,\sigma), -c, c\right)
          \label{eq:td3_target}
        \end{equation}
  \item \textbf{Delayed actor update}: the actor is updated every $d=2$
        critic iterations, reducing error propagation.
  \item \textbf{Target policy smoothing}: the noise $\epsilon$
        in~\eqref{eq:td3_target} regularises the target by smoothing peaks of the
        $Q$ function.
\end{enumerate}

\subsubsection{SAC — Soft Actor-Critic}

SAC~\parencite{Haarnoja2018} adopts a maximum-entropy optimisation framework.
The augmented objective is:
\begin{equation}
  J(\pi) = \E\!\left[\sum_t \gamma^t \Bigl(r(s_t, a_t)
         + \alpha\,\mathcal{H}\bigl(\pi(\cdot|s_t)\bigr)\Bigr)\right]
  \label{eq:sac_objective}
\end{equation}
where $\mathcal{H}(\pi(\cdot|s)) = -\E_{a\sim\pi}[\log\pi(a|s)]$ is the entropy and
$\alpha > 0$ is a temperature coefficient learned automatically:
\begin{equation}
  J(\alpha) = \E_{a\sim\pi}\!\left[-\alpha\log\pi(a|s) - \alpha\bar{\mathcal{H}}\right],
  \qquad \bar{\mathcal{H}} = -\dim(\mathcal{A})
  \label{eq:sac_alpha}
\end{equation}
The policy is parameterized as a squashed Gaussian through reparameterization:
\begin{equation}
  a_t = \tanh\!\left(\mu_\theta(s_t) + \varepsilon \odot \sigma_\theta(s_t)\right),
  \quad \varepsilon \sim \mathcal{N}(\mathbf{0}, \mathbf{I})
  \label{eq:sac_policy}
\end{equation}

\subsubsection{PPO --- Proximal Policy Optimization}

PPO~\parencite{Schulman2017} optimizes a clipped objective function to prevent
overly aggressive policy updates:
\begin{equation}
  \mathcal{L}^\text{CLIP}(\theta) =
  \E_t\!\left[\min\!\left(r_t(\theta)\,\hat{A}_t,\;
  \mathrm{clip}\!\left(r_t(\theta), 1-\varepsilon, 1+\varepsilon\right)\hat{A}_t\right)\right]
  \label{eq:ppo_clip}
\end{equation}
where $r_t(\theta) = \pi_\theta(a_t|s_t)/\pi_{\theta_\text{old}}(a_t|s_t)$ is the
probability ratio and $\hat{A}_t$ is the advantage estimated via GAE~\parencite{Schulman2016}:
\begin{equation}
  \hat{A}_t = \sum_{l=0}^{\infty} (\gamma\lambda)^l\,\delta_{t+l}, \qquad
  \delta_t = r_t + \gamma\,V(s_{t+1}) - V(s_t)
  \label{eq:gae}
\end{equation}
with $\lambda \in [0,1]$ controlling the bias--variance trade-off. PPO is an
\textit{on-policy} algorithm: it uses only recent trajectories, which implies lower
data efficiency but greater stability compared with DDPG.

\subsection{Applications of DRL to Musculoskeletal Control}

\textcite{Rajeswaran2017} use natural policy-gradient methods to learn dexterous
movements in musculoskeletal models, demonstrating the feasibility of \DRL{} on
highly redundant systems.

The \textit{Learning to Run} challenge~\parencite{Kidzinski2018} stimulated a
worldwide competition on the control of a bipedal musculoskeletal model, revealing
the superiority of \textit{off-policy} algorithms with replay memory for long-horizon
tasks. \textcite{Song2021} extend this approach to modeling human locomotion
control in a more general neuromechanical framework.

\textcite{Peng2018} introduce DeepMimic, which combines human motion-capture
data and \DRL{} to learn physically realistic behaviors, demonstrating the value of
combining \textit{imitation learning} with \DRL{}.

\textcite{Caggiano2022} present MyoSuite, with highly detailed hand and wrist
models ($>50$~muscles) for the study of fine motor control and rehabilitation, and
simultaneously launch MyoChallenge as a public benchmark.
\textcite{Codol2023} propose MotorNet, a differentiable simulator of a two-joint,
nine-muscle planar arm, fully integrated into PyTorch and enabling training by
backpropagation through physics.

\subsection{Domain Randomization and Sim-to-Real Transfer}

Domain Randomization~\parencite{Tobin2017, Peng2018SimToReal} consists of training
the agent on distributions of random physical parameters (masses, segment lengths,
muscle parameters, perturbations) in order to improve the robustness and
generalizability of the learned policies. This technique, initially developed for
sim-to-real transfer in robotics, applies naturally to musculoskeletal control by
making policies robust to inter-individual variability and modeling errors.

\subsection{Reproducibility and Statistical Rigor in \DRL{}}

\textcite{Henderson2018} highlighted the reproducibility fragility of published
\DRL{} results: inter-seed variance is often of the same order as the gap between
algorithms. \textcite{Agarwal2021} formalized a set of good practices (bootstrap
confidence intervals, \textit{interquartile mean}, performance profiles) that are now
considered standard for any serious benchmarking effort. Our experimental protocol
(Chapter~\ref{chap:impl}) adopts these recommendations.

% -----------------------------------------------------------
\section{Synthesis and Positioning}

The analysis of the state of the art reveals several gaps:

\begin{enumerate}[label=(\roman*)]
  \item Most studies compare \DRL{} algorithms on models without an explicit,
        mathematically complete Hill-type muscle model (SEE, PEE, pennation).
  \item The systematic analysis of emerging motor strategies, with quantitative
        comparison to human EMG data and synergy decomposition, remains underdeveloped.
  \item Rigorous comparative benchmarks including classical controller baselines
        (impedance PD) and adopting modern statistical standards~\parencite{Agarwal2021}
        are rare.
  \item The effect of Domain Randomization on the robustness of musculoskeletal
        policies, applied across all compared algorithms, has been little studied in a
        systematic way.
  \item The analysis of real-time computational feasibility (inference latency
        $<$~\SI{5}{\milli\second}) for clinical applications is rarely addressed.
  \item[\textit{(vi)}] \textbf{No study has directly benchmarked the per-step
        computational cost of physics-based muscle control} (Newton--Raphson
        equilibrium solving over all muscles) \textbf{against RL policy inference},
        making it impossible to quantify the computational argument for replacing the
        solver with a trained network. This gap leaves the core trade-off ---
        training cost vs.\ deployment efficiency --- empirically unresolved.
\end{enumerate}

Our contribution is distinguished by points~C1--C6 detailed in
Section~\ref{sec:contributions}, with~C6 specifically addressing gap~(vi).
